\section{Introduction}
\label{sec:intro}

Pauses, stalls and periods of inactivity in keystroke logs are among the most widely used process measures in
educational data mining. In introductory programming they serve to estimate time on task
\cite{hart2023time,hart2025time,edwards2022engagement} and are related to course performance
\cite{shrestha2022pausing,koutcheme2026keystroke}; in writing research they measure planning and fluency
\cite{roeser2024disfluencies,torrance2024methods,chukharev2014pauses,guo2020scenario}. Every such measure starts from a
definition: how long a gap must be, relative to what, over which stretch of the log, and what counts as producing text.
These choices differ from study to study. Time-on-task thresholds have ranged from five minutes to one hour
\cite{hart2025time}, and fixed pause thresholds in writing are acknowledged to have no theoretical ground between about
250~ms and 10~s \cite{torrance2024methods}.

That analytic choices can change conclusions is known. In learning analytics the strategy for estimating time on task
changed the results of the same regression \cite{kovanovic2015time}; operationalisations of wheel spinning disagreed on
most cases \cite{zhang2019wheel}; models of student success differ across a multiverse of reasonable choices
\cite{tang2025multiverse}; and in eye-movement research a multiverse of cleaning procedures left the existence of an
effect intact while changing its size \cite{godwin2025multiverse}. What has not been examined is how far the conclusions
that readers draw from a keystroke pause or stall label---its association with the quality of a product, comparisons of
rates between datasets, and the count and length structure within writers---move across definitions, and which of them
stay.

We answer this with a diagnostic. We take one window-based stall detector, which labels a stop in net character
production once a condition has held for a window and then a threshold (more than 15~s in our grid), and vary its four
choices (the window, the threshold, the condition for a stop and the segmentation of observation) over 89 definitions
around its default. Because the grid centres on one detector, we add a second detector of the kind used in writing
research, which counts inter-key intervals above a fixed threshold of 2 to 10~s (18 definitions). Both are applied to
four datasets released for research, from three sources: two editions of a CS1 course and two writing corpora. Each
conclusion is reported with an interval from resampling writers, so that the spread across definitions can be set
beside sampling variation. We do not evaluate the labels as predictors or as triggers for an intervention; we ask the
prior question of whether the conclusions that such uses rest on are stable.

The main result concerns the association with essay quality. By criteria set before the computations, an association
is present for most definitions, but its direction changes within each detector, and it changes again with
adjustment: given essay length, which correlates with the score at 0.75, no definition is negative
(Sections~\ref{sec:res-score} and~\ref{sec:res-adjust}). The same diagnostic applied to comparisons between datasets
and to the structure within writers shows that a rate ratio between the two CS1 editions reverses for a fifth of the
definitions and that overdispersion within writers depends on the definition, whereas the largest component of stall
length stays (Section~\ref{sec:res-other}).

We contribute a reporting convention for pause and stall rates in educational data mining: give the value of each
axis, the range of a conclusion across definitions and adjustments with intervals, and a three-valued reading of its
direction (consistent, negligible, or dependent on the definition). The score association in one corpus is the worked
example.
